\begin{table}[p]\centering
    \caption{\\Aggregated investment effects with heterogeneous marginal effects}\footnotesize
    \label{tabapp:agg_het}
    \begin{tabular}{lrrr}
    \toprule
    & 1933 & 1934 & After \\
    \midrule
    \multicolumn{4}{c}{\textit{Panel A: All firms}}\\
    \midrule
    Gold clause effect in \% (baseline)      & -1.77 & -1.29 & -0.21 \\
    Gold clause effect in \% (with rating)      & -1.41 & -1.79 & -0.51 \\
    Gold clause effect in \% (with size)      & -1.34 & -0.96 & -0.34 \\
    \midrule
    \multicolumn{4}{c}{\textit{Panel B: $\tilde{d}>0$ firms}}\\
    \midrule
    Gold clause effect in \% (baseline)      & -3.26 & -2.31 & -0.40 \\
    Gold clause effect in \% (with rating)      & -2.60 & -3.21 & -0.96 \\
    Gold clause effect in \% (with size)      & -2.47 & -1.72 & -0.65 \\
    \bottomrule

    \end{tabular}\\

    \vspace*{3mm} \justifying \noindent
\footnotesize{\textit{Notes.} This table extends the aggregation in Table \ref{tab:agg} by allowing for heterogeneous marginal effects of gold clause exposure across firm subgroups. Rows labeled ``with rating'' allow the marginal effect to differ between high-rated and low-rated firms (Ba or below in 1930), using coefficients from Table \ref{tab:credit_rating}. Rows labeled ``with size'' allow the marginal effect to differ between large and small firms (above and below median assets). The baseline row reproduces the uniform-effect estimates from Table \ref{tab:agg} for comparison. ``After'' denotes the post-resolution period (1935--1940). See Section \ref{sec:agg} for details.}
\end{table}